\documentclass[10pt,a4paper]{amsart}
\usepackage[margin=19mm]{geometry}
\usepackage{amsmath,amssymb,mathtools}
\usepackage[scr=rsfs,cal=euler]{mathalfa}
\usepackage{xfrac}
\usepackage[unicode,hidelinks,bookmarks=false]{hyperref}
\newcommand{\ie}{i.e.\ }
\newcommand{\cf}{cf.\ }
\newcommand{\resp}{respectively\ }
\newcommand{\Subsection}{\subsection}
\providecommand{\nobreakdash}{\nobreak\hbox{-}}
\setlength{\emergencystretch}{2em}
\multlinegap0pt
\usepackage{amssymb}
\usepackage{xcolor}
\newtheorem{theorem}{Theorem}
\title{Periodic Shadowing for Set-Valued Maps}
\author{M.C. Oliveira}
\date{Source: arXiv:2602.12399v2 (CC BY 4.0)}
\begin{document}
\maketitle

\noindent\emph{Source and licence.} Periodic Shadowing for Set-Valued Maps. Version arXiv:2602.12399v2 (CC BY 4.0), distributed by arXiv under the Creative Commons Attribution 4.0 International licence, http://creativecommons.org/licenses/by/4.0/, as recorded in the arXiv licence metadata for this exact version and shown on the abstract page https://arxiv.org/abs/2602.12399v2. Full licence notice: \texttt{licenses/arxiv-2602.12399-CC-BY-4.0.txt}.\par\smallskip

\noindent\textit{Editorial scope.} Counted unit: the article's theorem thm:involution (source0.tex lines 289--306). The article's complete original proof of it is delivered with it (source0.tex lines 308--368, one \texttt{proof} environment of 61 lines). No mathematics is written, repaired or re-derived: the statement and the proof are the article's own lines, in the article's own notation, and no retained line is substituted or re-flowed.  No further source-local context is needed: every cross-reference of every retained line resolves inside the two delivered spans.  Nothing was omitted from any retained span, so the package carries no omission record.  No retained line contains a citation, so the wrapper carries no bibliography and no bibliography entry.  No retained line contains a figure environment, an image-inclusion command or drawing code, so no image is delivered and the package declares no asset dependency.  Editorial formatting only, in addition: an AMS \texttt{amsart} preamble that restates the article's own declarations and macros used by the retained lines, namely the environments \texttt{theorem} on separate counters, together with the standard packages those lines need.  The retained environments print in this document as Theorem 1 in order of appearance, whereas the article prints the same items with the numbering of its own full document.  The complete native source of the article as distributed in the arXiv e-print for this exact version is bundled unchanged as \texttt{sources/194545/source0.tex}.
\begin{theorem}[Construction via Isometric Involutions]
\label{thm:involution}
Let $(X,d)$ be a metric space and let $f:X\to X$ be a map. Assume that $R:X\to X$ is an isometric involution satisfying $
R\circ f=f\circ R$.
Define the set-valued map $F:X\to 2^X$ by $F(x)=\{f(x),R(f(x))\}$.

Then the following assertions hold:

\begin{enumerate}
    \item If $f$ has the shadowing property, then $F$ has the shadowing property.

    \item If $f$ has the periodic shadowing property, then $F$ has the periodic shadowing property.

    \item If $f$ is topologically transitive, then $F$ is topologically transitive.

    \item If $f$ is chain transitive, then $F$ is chain transitive.
\end{enumerate}
\end{theorem}

\begin{proof}

We first prove (1). Let $\varepsilon>0$ and let $\delta>0$ be given by the shadowing property of $f$. Let $
\{x_n\}_{n\in\mathbb N}$ be a $\delta$-pseudo-orbit of $F$. We inductively construct a $\delta$-pseudo-orbit $
\{y_n\}_{n\in\mathbb N}$ of $f$.

Choose $y_0\in\{x_0,R(x_0)\}$. Assume first that $y_0=x_0$. Since $d(x_1,F(x_0))<\delta$, if $d(x_1,f(x_0))<\delta$, set $y_1=x_1$; otherwise define $
y_1=R(x_1)$. In this case,
\[
d(R(x_1),f(x_0))
=
d(x_1,R(f(x_0)))
<
\delta,
\]
since $R$ is an isometric involution commuting with $f$. Hence
$d(y_1,f(y_0))<\delta$. The case $y_0=R(x_0)$ is analogous.

Proceeding inductively, suppose that $y_n\in\{x_n,R(x_n)\}$ has been chosen so that $d(y_n,f(y_{n-1}))<\delta$. If $d(x_{n+1},f(x_n))<\delta$, define
$y_{n+1}=x_{n+1}$ when $y_n=x_n$, and $y_{n+1}=R(x_{n+1})$ otherwise. 

If instead $d(x_{n+1},R(f(x_n)))<\delta$, reverse the choice. Since $R$ is an isometric involution commuting with $f$, we obtain $d(y_{n+1},f(y_n))<\delta$ in either case. Thus, $\{y_n\}_{n\in\mathbb N}$ is a $\delta$-pseudo-orbit of $f$.

By the shadowing property of $f$, there exists $z\in X$ such that $
d(f^n(z),y_n)<\varepsilon$ for all $n\in\mathbb N$. 

Define $w_n=f^n(z)$ if $y_n=x_n$, and $w_n=R(f^n(z))$ if $y_n=R(x_n)$. Since
$R\circ f=f\circ R$, the sequence $\{w_n\}_{n\in\mathbb N}$ is an orbit of $F$, and 
\[
d(x_n,w_n)
=
d(y_n,f^n(z))
<
\varepsilon
\]
for all $n\in\mathbb N$, proving (1).

Assertion (2) follows similarly. Given a periodic $\delta$-pseudo-orbit of $F$, the above construction produces a periodic $\delta$-pseudo-orbit of $f$. By periodic shadowing, there exists a periodic orbit of $f$ shadowing it, which induces a periodic orbit of $F$.

Now, let $U,V\subset X$ be nonempty open sets. Since $f$ is topologically transitive, there exist $x\in U$ and $n\geq 0$ such that $f^n(x)\in V$.

We construct an orbit of $F$ by choosing $x_{k+1}=f(x_k)$ for every $k\geq 0$. Since $f(x_k)\in F(x_k)$, the sequence $\{x_k\}$ is an orbit of $F$.

Moreover, $x_0=x\in U$ and $x_n=f^n(x)\in V$. Hence $F$ is topologically transitive.

Consider $\delta>0$ and let $x,y\in X$. Since $f$ is chain transitive, there exists a $\delta$-chain
\[
x=x_0,x_1,\ldots,x_n=y
\]
such that $d(f(x_i),x_{i+1})<\delta$ for all $0\leq i\leq n-1$. Since $f(x_i)\in F(x_i)$, we obtain
\[
d(F(x_i),x_{i+1})
\leq d(f(x_i),x_{i+1})
<\delta.
\]
Thus,
\[
x=x_0,x_1,\ldots,x_n=y
\]
is also a $\delta$-chain for $F$. Therefore, $F$ is chain transitive.
\end{proof}

\end{document}
